\SourceLeaf{417}


\Needspace{9\baselineskip}\begin{center}{\Large SIEBENTER ABSCHNITT}\par\medskip ÜBER\par\medskip {\large DIE GLEICHUNGEN, WELCHE DIE KREISTEILUNGEN BESTIMMEN.}\end{center}
\BookRule{0.5\textwidth}

\Needspace{6\baselineskip}\Article{335}
Unter den glänzendsten Fortschritten, welche die Mathematik durch die Arbeiten neuerer Mathematiker erfahren hat, nimmt die Theorie der vom Kreis abhängigen Funktionen zweifellos einen besonders hervorragenden Platz ein. Die größten Mathematiker neuerer Zeit haben dieser wunderbaren Art von Größen, zu der wir bei den verschiedenartigsten Untersuchungen so häufig gelangen und deren Hilfe kein Teil der gesamten Mathematik entbehren kann, ihren Fleiß und Scharfsinn so beharrlich gewidmet und daraus ein so weitreichendes Gebiet geschaffen, dass man kaum hätte erwarten können, irgendein Teil dieser Theorie, geschweige denn ein elementarer und gewissermaßen an ihrer Schwelle liegender, könne noch wesentliche Fortschritte erfahren. Ich spreche von der Theorie der trigonometrischen Funktionen zu den mit dem Umfang kommensurablen Bögen, oder von der Theorie der regelmäßigen Vielecke; dieser Abschnitt wird zeigen, wie wenig davon bisher entwickelt worden ist. Die Leser könnten sich wundern, dass eine solche Untersuchung gerade in diesem Werk angestellt wird, das hauptsächlich einem auf den ersten Blick ganz andersartigen Gegenstand gewidmet ist; doch wird die Behandlung selbst hinreichend deutlich machen, wie eng dieser Gegenstand mit der höheren Arithmetik verbunden ist.

Die Grundsätze der Theorie, deren Darstellung wir beginnen, reichen übrigens viel weiter, als wir sie hier ausdehnen. Sie lassen sich nicht nur auf die Kreisfunktionen, sondern mit gleichem\SourceLeaf{418}
Erfolg auf viele andere transzendente Funktionen anwenden, zum Beispiel auf diejenigen, die vom Integral \(\int\frac{dx}{\surd{(1-x^4)}}\) abhängen, und außerdem auf verschiedene Arten von Kongruenzen. Da wir aber über jene transzendenten Funktionen ein umfangreiches besonderes Werk vorbereiten und die Kongruenzen in der Fortsetzung der arithmetischen Untersuchungen ausführlich behandelt werden sollen, schien es uns angebracht, hier nur die Kreisfunktionen zu betrachten. Ja, auch diese, die wir in größter Allgemeinheit behandeln könnten, werden wir mit den in den folgenden Artikeln darzulegenden Hilfsmitteln auf den einfachsten Fall zurückführen, sowohl der Kürze wegen als auch, damit die völlig neuen Grundsätze dieser Theorie leichter verstanden werden.

\Needspace{9\baselineskip}\begin{center}\textit{Die Untersuchung wird auf den einfachsten Fall zurückgeführt, in dem die Anzahl der Teile, in welche der Kreis geteilt werden soll,} \textit{eine Primzahl ist.}\end{center}


\Needspace{6\baselineskip}\Article{336}
Bezeichnet man den Umfang des Kreises oder vier rechte Winkel mit \(P\) und setzt voraus, dass \(m,\allowbreak n\) ganze Zahlen sind und \(n\) das Produkt aus den paarweise teilerfremden Faktoren \(a,\allowbreak b,\allowbreak c\) usw. ist, so lässt sich der Winkel \(A=\frac{mP}{n}\) nach Artikel 310 auf die Form \(A=(\frac\alpha a+\frac\beta b+\frac\gamma c+\mathrm{usw.})P\) bringen; die ihm entsprechenden trigonometrischen Funktionen lassen sich nach bekannten Verfahren aus den Funktionen zu den Teilen \(\frac{\alpha P}a,\allowbreak \frac{\beta P}b\) usw. ableiten. Da man für \(a,\allowbreak b,\allowbreak c\) usw. Primzahlen oder Potenzen von Primzahlen nehmen kann, genügt es offenbar, die Teilung des Kreises in Teile zu betrachten, deren Anzahl eine Primzahl oder eine Primzahlpotenz ist; das Vieleck mit \(n\) Seiten erhält man dann unmittelbar aus den Vielecken mit \(a,\allowbreak b,\allowbreak c\) usw. Seiten. Hier werden wir die Untersuchung jedoch auf den Fall beschränken, in dem die Anzahl der Teile, in welche der Kreis geteilt werden soll, eine (ungerade) Primzahl ist, und zwar vor allem aus folgendem Grund. Bekanntlich lassen sich die Kreisfunktionen zum Winkel \(\frac{mP}{pp}\) aus den Funktionen zu \(\frac{mP}p\) durch Lösung einer Gleichung vom Grad \(p\) ableiten, und ebenso aus jenen durch eine Gleichung vom selben Grad die Funktionen zu \(\frac{mP}{p^3}\) usw.; ist also das Vieleck mit \(p\) Seiten schon gegeben, so erfordert die Bestimmung des Vielecks mit \(p^\lambda\) Seiten notwendig die Lösung von \(\lambda-1\) Gleichungen vom Grad \(p\). Zwar lässt sich die folgende Theorie auch auf diesen Fall ausdehnen, doch gelangt man auf diesem Weg ebenfalls zu ebenso vielen Gleichungen vom Grad \(p\), die sich, wenn \(p\) eine Primzahl ist, in keiner Weise auf niedrigere Grade zurückführen lassen. So wird zum Beispiel unten gezeigt, dass sich das Vieleck mit \(17\) Seiten geometrisch konstruieren lässt; bei der Bestimmung des Vielecks mit \(289\) Seiten kann man dagegen eine Gleichung vom Grad \(17\) keinesfalls vermeiden.
\SourceLeaf{419}

\Needspace{9\baselineskip}\begin{center}\textit{Gleichungen für die trigonometrischen Funktionen der Bögen, die ein Teil oder mehrere Teile des ganzen Umfangs sind;} \textit{Zurückführung der trigonometrischen Funktionen auf die Wurzeln der Gleichung \(x^n-1=0\).}\end{center}


\Needspace{6\baselineskip}\Article{337}
Es ist bekannt, dass sich die trigonometrischen Funktionen aller Winkel \(\frac{kP}n\), wobei \(k\) sämtliche Zahlen \(0,\allowbreak 1,\allowbreak 2\ldots n-1\) durchläuft, durch die Wurzeln von Gleichungen vom Grad \(n\) ausdrücken lassen: nämlich die \textit{Sinus} durch die Wurzeln von (I) \[\begin{aligned}
x^n-\frac14nx^{n-2}+\frac1{16}\frac{n.n-3}{1.2}x^{n-4}
&-\frac1{64}\frac{n.n-4.n-5}{1.2.3}x^{n-6}\\
&+\mathrm{usw.}\pm\frac1{2^{n-1}}nx=0
\end{aligned}\] die \textit{Kosinus} durch die Wurzeln von (II) \[\begin{aligned}
x^n-\frac14nx^{n-2}+\frac1{16}\frac{n.n-3}{1.2}x^{n-4}
&-\frac1{64}\frac{n.n-4.n-5}{1.2.3}x^{n-6}\\
&+\mathrm{usw.}\pm\frac1{2^{n-1}}nx-\frac1{2^{n-1}}=0
\end{aligned}\] und schließlich die \textit{Tangenten} durch die Wurzeln von (III) \[\begin{aligned}
x^n-\frac{n.n-1}{1.2}x^{n-2}&+\frac{n.n-1.n-2.n-3}{1.2.3.4}x^{n-4}\\
&-\mathrm{usw.}\pm nx=0
\end{aligned}\] Diese Gleichungen (die allgemein für jeden ungeraden Wert von \(n\) gelten, II aber auch für gerade Werte) lassen sich durch Setzen von \(n=2m+1\) leicht auf den Grad \(m\) zurückführen: I und III, indem man die linke Seite durch \(x\) dividiert und \(y\) für \(xx\) einsetzt. Gleichung II hat offenbar die Wurzel \(x=1\ (=\cos0)\), und von den übrigen sind stets je zwei gleich (\(\cos\frac Pn=\cos\frac{(n-1)P}n,\allowbreak \ \cos\frac{2P}n=\cos\frac{(n-2)P}n\) usw.); daher ist ihre linke Seite durch \(x-1\) teilbar und der Quotient ein Quadrat. Zieht man daraus die Quadratwurzel, so wird Gleichung II auf \[\begin{aligned}
x^m+\frac12x^{m-1}-\frac14(m-1)x^{m-2}&-\frac18(m-2)x^{m-3}\\
&+\frac1{16}\frac{m-2.m-3}{1.2}x^{m-4}\\
&+\frac1{32}\frac{m-3.m-4}{1.2}x^{m-5}-\mathrm{usw.}=0
\end{aligned}\] zurückgeführt, deren Wurzeln die Kosinus der Winkel \(\frac Pn,\allowbreak \frac{2P}n,\allowbreak \frac{3P}n\ldots\frac{mP}n\) sind. Weitere Zurückführungen dieser Gleichungen, jedenfalls für den Fall, dass \(n\) eine Primzahl ist, standen bisher nicht zur Verfügung.

Keine dieser Gleichungen ist jedoch so leicht zu behandeln und für unseren Zweck so geeignet wie \(x^n-1=0\), deren Wurzeln bekanntlich sehr eng mit den Wurzeln jener zusammenhängen. Schreibt man nämlich der Kürze wegen \(i\) für die imaginäre Größe \(\surd{-1}\), so werden die Wurzeln der Gleichung \(x^n-1=0\) dargestellt durch \[\cos\frac{kP}n+i\sin\frac{kP}n=r,\]
\SourceLeaf{420}
wobei für \(k\) sämtliche Zahlen \(0,\allowbreak 1,\allowbreak 2\ldots n-1\) zu nehmen sind. Da \(\frac1r=\cos\frac{kP}n-i\sin\frac{kP}n\) gilt, werden also die Wurzeln der Gleichung I durch \(\frac1{2i}(r-\frac1r)\) oder durch \(i\frac{1-rr}{2r}\) dargestellt, die Wurzeln der Gleichung II durch \(\frac12(r+\frac1r)=\frac{1+rr}{2r}\) und schließlich die Wurzeln der Gleichung III durch \(\frac{i(1-rr)}{1+rr}\). Deshalb werden wir unsere Untersuchung auf die Betrachtung der Gleichung \(x^n-1=0\) gründen und dabei \(n\) als ungerade Primzahl voraussetzen. Um die Reihenfolge der Untersuchungen nicht unterbrechen zu müssen, schicken wir hier den folgenden Hilfssatz voraus.

\clearpage
\Needspace{6\baselineskip}\Article{338}
\textsc{Aufgabe.} \textit{Zur gegebenen Gleichung} \[(W)\ldots z^m+Az^{m-1}+\mathrm{usw.}=0,\] \textit{ist eine Gleichung \((W')\) zu finden, deren Wurzeln die \(\lambda\)-ten Potenzen der Wurzeln der Gleichung \((W)\) sind,} \textit{wobei \(\lambda\) einen gegebenen positiven ganzzahligen Exponenten bezeichnet.} \textit{Lösung.} Bezeichnet man die Wurzeln der Gleichung \(W\) mit \(a,\allowbreak b,\allowbreak c\) usw., so müssen die Wurzeln der Gleichung \(W'\) gleich \(a^\lambda,\allowbreak b^\lambda,\allowbreak c^\lambda\) usw. sein. Nach dem bekannten Newtonschen Satz kann man aus den Koeffizienten der Gleichung \(W\) die Summen beliebiger Potenzen der Wurzeln \(a,\allowbreak b,\allowbreak c\) usw. finden. Man suche also die Summen \[\begin{gathered}
a^\lambda+b^\lambda+c^\lambda+\mathrm{usw.},\quad
a^{2\lambda}+b^{2\lambda}+c^{2\lambda}\ \mathrm{usw.}\ \mathrm{usw.}\\
\text{bis zu}\quad a^{m\lambda}+b^{m\lambda}+c^{m\lambda}+\mathrm{usw.}
\end{gathered}\] aus denen man dann umgekehrt nach demselben Satz die Koeffizienten der Gleichung \(W'\) ableiten kann. \textit{Q. E. F.} — Zugleich geht daraus hervor, dass, wenn alle Koeffizienten in \(W\) rational sind, auch alle in \(W'\) rational werden. Auf einem anderen Weg lässt sich zwar beweisen, dass, wenn jene alle ganze Zahlen sind, auch diese alle ganze Zahlen werden; bei diesem für unseren Zweck nicht so nötigen Satz halten wir uns hier jedoch nicht auf.

\Needspace{6\baselineskip}\Article{339}
Die Gleichung \(x^n-1=0\) (unter der von nun an stets mitzuverstehenden Voraussetzung, dass \(n\) eine ungerade Primzahl ist) hat nur eine reelle Wurzel, \(x=1\); die übrigen \(n-1\) Wurzeln sind in der Gleichung \[x^{n-1}+x^{n-2}+\mathrm{usw.}+x+1=0\] enthalten und alle imaginär. Ihre Menge werden wir mit \(\Omega\) und die Funktion \[x^{n-1}+x^{n-2}+\mathrm{usw.}+x+1\quad\text{mit}\quad X\]
\SourceLeaf{421}
bezeichnen. Ist also \(r\) irgendeine Wurzel aus \(\Omega\), so gilt \(1=r^n=r^{2n}\) usw. und allgemein \(r^{en}=1\) für jeden ganzzahligen Wert von \(e\), ob positiv oder negativ; daraus ist klar, dass, wenn \(\lambda,\allowbreak \mu\) nach \(n\) kongruente ganze Zahlen sind, \(r^\lambda=r^\mu\) gilt. Sind dagegen \(\lambda,\allowbreak \mu\) nach dem Modul \(n\) inkongruent, so sind \(r^\lambda\) und \(r^\mu\) ungleich; denn in diesem Fall lässt sich eine ganze Zahl \(\nu\) so wählen, dass \((\lambda-\mu)\nu\equiv1\ (\mathrm{mod.}\ n)\) gilt, woraus \(r^{(\lambda-\mu)\nu}=r\) folgt und daher \(r^{\lambda-\mu}\) sicher nicht \(=1\) ist. Ferner ist klar, dass jede Potenz von \(r\) ebenfalls eine Wurzel der Gleichung \(x^n-1=0\) ist; da die Größen \(1\ (=r^0),\allowbreak r,\allowbreak rr\ldots r^{n-1}\) alle verschieden sind, stellen sie also sämtliche Wurzeln der Gleichung \(x^n-1=0\) dar, und somit stimmen \(r,\allowbreak rr,\allowbreak r^3\ldots r^{n-1}\) mit \(\Omega\) überein. Allgemeiner folgt daraus leicht, dass \(\Omega\) mit \(r^e,\allowbreak r^{2e},\allowbreak r^{3e}\ldots r^{(n-1)e}\) übereinstimmt, wenn \(e\) irgendeine durch \(n\) nicht teilbare ganze Zahl ist, ob positiv oder negativ. Es gilt daher \[X=(x-r^e)(x-r^{2e})(x-r^{3e})\ldots(x-r^{(n-1)e})\] woraus \[\begin{gathered}
r^e+r^{2e}+r^{3e}+\ldots+r^{(n-1)e}=-1,\\
\text{und}\quad1+r^e+r^{2e}+\ldots+r^{(n-1)e}=0
\end{gathered}\] folgt.

Zwei Wurzeln wie \(r\) und \(\frac1r\ (=r^{n-1})\) oder allgemein \(r^e\) und \(r^{-e}\) werden wir \textit{reziprok} nennen; offenbar ist das Produkt der beiden einfachen Faktoren \(x-r\) und \(x-\frac1r\) reell, \(=xx-2x\cos\omega+1\), wobei der Winkel \(\omega\) entweder dem Winkel \(\frac Pn\) oder einem seiner Vielfachen gleich ist.

\Needspace{6\baselineskip}\Article{340}
Bezeichnet man also eine Wurzel aus \(\Omega\) mit \(r\), so werden sämtliche Wurzeln der Gleichung \(x^n-1=0\) durch Potenzen von \(r\) ausgedrückt; ein auf beliebige Weise aus mehreren Wurzeln dieser Gleichung gebildetes Produkt lässt sich daher durch \(r^\lambda\) darstellen, wobei \(\lambda\) entweder \(0\) oder positiv und \(<n\) ist. Bezeichnet nun \(\varphi(t,\allowbreak u,\allowbreak v\ldots)\) eine ganzrationale algebraische Funktion der Unbestimmten \(t,\allowbreak u,\allowbreak v\) usw., die sich als Summe von Gliedern der Form \(ht^\alpha u^\beta v^\gamma\ldots\) ausdrücken lässt, so ist klar: Setzt man für \(t,\allowbreak u,\allowbreak v\) usw. einige der Wurzeln der Gleichung \(x^n-1=0\) ein, etwa \(t=a,\allowbreak u=b,\allowbreak v=c\) usw., so lässt sich \(\varphi(a,\allowbreak b,\allowbreak c\ldots)\) auf die Form \[A+A^{\prime} r+A^{\prime\prime} rr+A^{\prime\prime\prime} r^3+\ldots+A^\nu r^{n-1}\] bringen, wobei die Koeffizienten \(A,\allowbreak A^{\prime}\) usw. (von denen auch einige fehlen und somit \(=0\) sein können) bestimmte Größen sind. Außerdem sind alle diese Koeffizienten ganze Zahlen, wenn sämtliche bestimmten Koeffizienten in \(\varphi(t,\allowbreak u,\allowbreak v\ldots)\), das heißt alle
\SourceLeaf{422}
\(h\), ganze Zahlen sind. Setzt man danach für \(t,\allowbreak u,\allowbreak v\ldots\) jeweils \(aa,\allowbreak bb,\allowbreak cc\ldots\) ein, so wird jedes Glied wie \(ht^\alpha u^\beta v^\gamma\ldots\), das sich vorher auf \(r^\sigma\) zurückführen ließ, nun zu \(r^{2\sigma}\). Daraus folgt leicht \[\varphi(aa,bb,cc\ldots)=A+A^{\prime} rr+A^{\prime\prime} r^4+A^{\prime\prime\prime} r^6+\ldots+A^\nu r^{2n-2}\] Ebenso gilt allgemein für jeden ganzzahligen Wert von \(\lambda\) \[\varphi(a^\lambda,b^\lambda,c^\lambda\ldots)=A+A^{\prime} r^\lambda+A^{\prime\prime} r^{2\lambda}+\ldots+A^\nu r^{(n-1)\lambda}\] Dieser Satz ist von größter Bedeutung und bildet die Grundlage der folgenden Untersuchungen. — Daraus folgt außerdem \[\varphi(1,1,1\ldots)=\varphi(a^n,b^n,c^n\ldots)=A+A^{\prime}+A^{\prime\prime}+\ldots+A^\nu\] sowie \[\begin{aligned}
\varphi(a,b,c\ldots)+\varphi(aa,bb,cc\ldots)&+\varphi(a^3,b^3,c^3\ldots)\\
&+\ldots+\varphi(a^n,b^n,c^n\ldots)=nA
\end{aligned}\] Diese Summe ist also immer eine ganze, durch \(n\) teilbare Zahl, wenn alle bestimmten Koeffizienten in \(\varphi(t,\allowbreak u,\allowbreak v\ldots)\) ganze Zahlen sind.

